\documentclass[12pt]{amsart}
%prepared in AMSLaTeX, under LaTeX2e
\addtolength{\oddsidemargin}{-1.1in} 
\addtolength{\evensidemargin}{-1.1in}
\addtolength{\topmargin}{-.9in}
\addtolength{\textwidth}{1.75in}
\addtolength{\textheight}{1.75in}

\renewcommand{\baselinestretch}{1.05}

\usepackage{verbatim} % for "comment" environment

\usepackage{palatino}

\usepackage[final]{graphicx}

\usepackage{tikz}
\usetikzlibrary{positioning}

\usepackage{bm,enumitem,xspace,fancyvrb}

\newtheorem*{thm}{Theorem}
\newtheorem*{defn}{Definition}
\newtheorem*{example}{Example}
\newtheorem*{problem}{Problem}
\newtheorem*{remark}{Remark}

\DefineVerbatimEnvironment{mVerb}{Verbatim}{numbersep=2mm,frame=lines,framerule=0.1mm,framesep=2mm,xleftmargin=4mm,fontsize=\footnotesize}

% macros
\usepackage{amssymb}
\newcommand{\bA}{\mathbf{A}}
\newcommand{\bB}{\mathbf{B}}
\newcommand{\bE}{\mathbf{E}}
\newcommand{\bF}{\mathbf{F}}
\newcommand{\bJ}{\mathbf{J}}

\newcommand{\bb}{\mathbf{b}}
\newcommand{\bc}{\mathbf{c}}
\newcommand{\bd}{\mathbf{d}}
\newcommand{\bi}{\mathbf{i}}
\newcommand{\bj}{\mathbf{j}}
\newcommand{\bk}{\mathbf{k}}
\newcommand{\br}{\mathbf{r}}
\newcommand{\bv}{\mathbf{v}}
\newcommand{\bw}{\mathbf{w}}
\newcommand{\bx}{\mathbf{x}}

\newcommand{\bzero}{\bm{0}}

\newcommand{\eps}{\epsilon}
\newcommand{\grad}{\nabla}
\newcommand{\ip}[2]{\ensuremath{\left<#1,#2\right>}}
\newcommand{\lam}{\lambda}
\newcommand{\lap}{\triangle}

\newcommand{\Null}{\operatorname{null}}
\newcommand{\rank}{\operatorname{rank}}
\newcommand{\range}{\operatorname{range}}
\newcommand{\trace}{\operatorname{tr}}

\newcommand{\RR}{\mathbf{R}}

\newcommand{\ZZ}{\mathbb{Z}}

\newcommand{\prob}[1]{\bigskip\noindent\textbf{#1.}\quad }
\newcommand{\exer}[2]{\prob{Exercise #2 on page #1}}

\newcommand{\pts}[1]{(\emph{#1 pts}) }
\newcommand{\epart}[1]{\bigskip\noindent\textbf{(#1)}\quad }
\newcommand{\ppart}[1]{\,\textbf{(#1)}\quad }

\newcommand{\Matlab}{\textsc{Matlab}\xspace}
\newcommand{\ds}{\displaystyle}

\def\vectwo#1#2{\begin{bmatrix}#1\\#2\end{bmatrix}}
\def\vecthree#1#2#3{\begin{bmatrix}#1\\#2\\#3\end{bmatrix}}
\def\vecfour#1#2#3#4{\begin{bmatrix}#1\\#2\\#3\\#4\end{bmatrix}}
\def\bbm{\begin{bmatrix}}
\def\ebm{\end{bmatrix}}


\begin{document}
\scriptsize \noindent Math 314 Linear Algebra \hfill Section 4.1 Example and Principles 
\normalsize\medskip


\thispagestyle{empty}
\begin{enumerate}
\item Find a basis for the four subspaces of $A.$ \\

$A =  \begin{bmatrix} 1 & 2 & 0 &0&-2\\ 
1 & 2 & -2&0&0 \\ 1 & 2 & -1&1&-4 \\2&4&3&-3&2\end{bmatrix} 
\quad \to \quad R = \begin{bmatrix} 1 &2&0& 0 & -2 \\0& 0 & 1 & 0&-1 \\ 0 & 0 & 0&1&-3 \\ 0&0&0&0&0  \end{bmatrix}.$\\

$A^{T} \quad \to \quad R = \begin{bmatrix} 1 &0& 0 & 5 \\0 & 1 & 0&0 \\ 0 & 0 &1&-3 \\ 0&0&0&0 \\ 0&0&0&0  \end{bmatrix}$\\
\vfill
\item Why are the principles below true?
	\begin{enumerate}
	\item If $\bv$ and $\bw$ are orthogonal, then $ \left\lVert \bv + \bw \right\rVert^2=\left\lVert \bv \right\rVert^2+ \left\lVert\bw \right\rVert^2.$
	\vspace{.75in}
	\item If $\bv$ and $\bw$ are orthogonal, then $\bv$ and $\bw$ are linearly independent.
	\vspace{.75in}
\newpage

	\item For any matrix $A$, $C(A^T) \perp N(A).$ 
	\vfill
	
	\vfill
	\item For any matrix $A$, $C(A) \perp N(A^T).$
	\vspace{.5in}
	\item If vector space $V$ has two bases $B_1=\{v_1,v_2,\cdots,v_m\}$ and $B_2 = \{w_1,w_2,\cdots, w_n\}$, then $m=n.$ (That is, the notion of the dimension of a vector space is well-defined.)
	\vfill
	\vfill
	\item  If the dimension of a vector space, $V$, is $d,$ then you can tell that a set of $d$ vectors is a basis of $V$ by checking whether it is...
\vfill
	\end{enumerate}
\end{enumerate}
\end{document}

